% Source: https://artofproblemsolving.com/community/c3h3163505_2018_fall_lmt_individual_round__lexington_mathematical_tournament
% Stable identity: aops-topic-3163505-q16
Use $C=(0,0),A=(3,0),B=(0,4)$. The circumcenter is $(\dfrac{3}{2},2)$. Perpendicularity to the radii gives tangent equations $3x-4y=9$ at $A$ and $-3x+4y=16$ at $B$. Hence $D=(0,-\dfrac{9}{4})$ and $E=(-\dfrac{16}{3},0)$. Line $DE$ is $27x+64y=-144$, while $AB$ is $4x+3y=12$. Solving gives $F=(\dfrac{48}{7},-\dfrac{36}{7})$. Therefore $BF=\sqrt{(\dfrac{48}{7})^2+(\dfrac{64}{7})^2}=\boxed{\dfrac{80}{7}}$.
\begin{asy}
usepackage("amsmath");
size(7cm);
pair A=(3,0),Bp=(0,4),C=(0,0),D=(0,-9/4),Ep=(-16/3,0),F=(48/7,-36/7);draw(circle((1.5,2),2.5));draw(A--Bp--C--cycle);draw(Ep--F);draw(A--D);draw(Bp--Ep);draw(Bp--F,dashed);draw(D--C,dashed);label("$A$",A,NE);label("$B$",Bp,N);label("$C$",C,W);label("$D$",D,W);label("$E$",Ep,W);label("$F$",F,SE);
\end{asy}
